\chapter*{Recomendaciones}
\addcontentsline{toc}{chapter}{Recomendaciones}

\begin{enumerate}
  \item Se recomienda que las instituciones bancarias que adopten la
        metodología propuesta realicen una calibración de los pesos de los
        criterios técnicos y regulatorios a su contexto específico antes de
        iniciar la Fase 2. Los pesos presentados en este trabajo
        (técnica 60\% / regulatoria 40\%) corresponden a una institución
        bancaria mediana peruana y pueden requerir ajuste según el tamaño
        de la organización, el marco regulatorio nacional y el nivel de
        adopción de herramientas de IA generativa en el equipo de
        desarrollo. La Técnica Delphi con expertos internos --desarrolladores
        senior, arquitectos y oficiales de cumplimiento-- es el instrumento
        recomendado para esta calibración.

  \item Dado el ritmo acelerado de evolución tanto del ecosistema de
        interfaz de usuario como de las herramientas de generación de
        código asistida por IA --con lanzamientos de versiones mayores
        anuales en Angular y React, y mejoras semestrales en asistentes
        como GitHub Copilot y Cursor-- se recomienda revisar los criterios
        técnicos de evaluación, y en particular el criterio de calidad del
        código generado por IA, con una periodicidad mínima de 12 meses. La
        metodología propuesta es compatible con revisiones parciales que
        actualicen únicamente los datos de la Fase 2 sin repetir las fases
        1, 3 y 4 completas.

  \item Se recomienda incorporar el Reglamento DORA (\emph{Digital
        Operational Resilience Act}), vigente en la Unión Europea desde enero
        de 2025, como criterio de filtro previo obligatorio para
        instituciones bancarias con presencia o aspiraciones en el mercado
        europeo, integrándolo junto con el criterio de gobernanza de IA
        propuesto en el Capítulo 3. Los requisitos de resiliencia
        operacional digital de DORA --gestión de riesgos TIC de terceros,
        pruebas de penetración basadas en amenazas (TLPT), notificación de
        incidentes-- imponen restricciones arquitectónicas específicas que
        deben verificarse antes de iniciar la evaluación comparativa de
        frameworks, ya que pueden eliminar opciones con dependencias en
        proveedores no conformes, incluidos proveedores de asistentes de IA.

  \item Se recomienda que las instituciones bancarias establezcan una
        política formal de gobernanza para el uso de asistentes de código
        basados en IA en el desarrollo de software bancario --con revisión
        humana obligatoria (\emph{human-in-the-loop}), trazabilidad de la
        procedencia del código sugerido y verificación de licenciamiento--
        de forma previa e independiente a la selección del framework de
        interfaz de usuario, dado que el criterio de gobernanza de IA
        propuesto en esta metodología presupone la existencia de dicha
        política institucional.

  \item Para investigaciones futuras, se recomienda extender la metodología
        propuesta a la selección de tecnologías de interfaz de usuario para
        canales móviles nativos e híbridos (React Native, Flutter, Ionic),
        donde las instituciones bancarias enfrentan decisiones tecnológicas
        igualmente críticas pero con un conjunto diferente de restricciones
        regulatorias y de calidad de código asistido por IA. La estructura
        metodológica de cuatro fases es directamente adaptable a este
        dominio con ajustes en los criterios de la Fase 2.

  \item Se recomienda que las instituciones bancarias que implementen la
        metodología documenten sistemáticamente sus lecciones aprendidas en
        un repositorio institucional de decisiones tecnológicas
        (\emph{Architecture Decision Records}), incluyendo los scorecard de
        evaluación completos, las actas de deliberación y los prompts o
        contextos utilizados para generar código durante las pruebas de
        concepto. Esta base de conocimiento reduciría el tiempo de ejecución
        en selecciones tecnológicas futuras y proporcionaría evidencia
        histórica para auditorías regulatorias de gobernanza tecnológica y
        de IA.
\end{enumerate}
